\chapter{Siemens S7-1500 and TIA Portal}

\section{S7-1500 PLC Architecture}

\subsection{Overview and Scan Cycle}

The Siemens S7-1500 is a high-performance PLC series designed for demanding automation
tasks requiring high processing speed, integrated motion control, and native PROFINET
communication. The installation at ICAM uses an S7-1500 CPU at IP address 192.168.1.10.

The S7-1500 CPU executes a deterministic scan cycle comprising three phases:

\begin{enumerate}[leftmargin=*, itemsep=3pt]
  \item \textbf{Input update:} Physical inputs from the I/O modules are read into the
    Process Image Input (PII), creating a synchronised snapshot of all \plcaddr{I} and
    \plcaddr{IW} values.
  \item \textbf{Program execution:} The user program (OB1 and called FCs/FBs) runs,
    reading from PII and writing results to the Process Image Output (PIO) and memory
    areas (\plcaddr{M}, \plcaddr{DB}).
  \item \textbf{Output update:} PIO values are transferred to the physical output modules
    (\plcaddr{Q}, \plcaddr{QW}).
\end{enumerate}

This three-phase structure guarantees that all logic within one scan cycle operates on a
consistent view of the I/O, preventing race conditions between input reading and output
writing. The scan cycle time is typically a few milliseconds --- much shorter than the
3-second poll rate of the web application.

\subsection{Memory Organisation}

Table~\ref{tab:plcmem} summarises the memory areas relevant to this project.

\begin{table}[H]
  \centering
  \caption{S7-1500 memory areas used in this project.}
  \label{tab:plcmem}
  \begin{tabular}{@{}L{2.8cm}L{2.0cm}L{3cm}L{5.2cm}@{}}
    \toprule
    \textbf{Area} & \textbf{Symbol} & \textbf{Access} & \textbf{Purpose} \\
    \midrule
    Process Image Input  & \plcaddr{I}/\plcaddr{IW} & Read-only (scan) & Snapshot of physical inputs \\
    Process Image Output & \plcaddr{Q}/\plcaddr{QW} & Written by OB1   & Drives physical outputs \\
    Bit Memory (Merkers) & \plcaddr{M}/\plcaddr{MD} & Read/write       & Global variables, scaled values \\
    Data Blocks          & \plcaddr{DB}             & Read/write       & Structured storage (Snap7 access) \\
    Timers / Counters    & \plcaddr{T}/\plcaddr{C}  & System-managed   & Timing and counting \\
    \bottomrule
  \end{tabular}
\end{table}

The deliberate use of \plcaddr{MD} (Merker Double-word, 32-bit REAL) for scaled
engineering-unit values decouples the scaling logic from the communication layer:
both OB1 and external tools (Snap7, Open Pipe) access the same scaled variables.

\section{TIA Portal V18 Configuration}

TIA Portal (Totally Integrated Automation Portal) is the unified Siemens engineering
environment for programming and configuring S7 PLCs, WinCC HMI, and drive systems.
The project uses TIA Portal V18.

\subsection{Hardware Configuration}

The hardware configuration defines the physical arrangement of CPU and I/O modules and
assigns \plcaddr{IW}/\plcaddr{QW} addresses to each module. The addresses depend directly
on the physical slot position of each module: changing the slot order changes all downstream
addresses. The hardware configuration is therefore the authoritative source for address
allocation.

For the thermal rig, analog input modules (AI) occupy specific slots, determining the
\plcaddr{IW} addresses listed in Table~\ref{tab:tags}. The analog output module determines
the \plcaddr{QW8} and \plcaddr{QW4} addresses for \code{Valve\_F1} and
\code{Valve\_F2}.

\subsection{Tag Table}

The TIA Portal tag table is the central register mapping symbolic names to memory addresses.
This project defines 21 process tags (Table~\ref{tab:tags}), retained after a cleanup
that removed 25 intermediate, redundant, or test variables from an initial set of 46.

The naming convention adopted uses English names with the suffix \code{\_HE1} for heat
exchanger variables (e.g., \code{Tin1\_HE1}, \code{Tout2\_HE1}, \code{Valve\_F1}).
Raw ADC signals carry the suffix \code{\_raw} or \code{\_skal}.

\subsection{OB1 and Scaling Function}

OB1 (Organisation Block~1) is the main cyclic program, called every scan cycle. It contains:

\begin{enumerate}[leftmargin=*, itemsep=3pt]
  \item \textbf{Scaling FC calls.} For each analog input channel, a Function block reads
    the raw \plcaddr{IW} value and computes the engineering-unit value using
    equation~\eqref{eq:scaling}, writing the result to the corresponding \plcaddr{MD}
    address. For example, raw \plcaddr{IW112} is scaled to \si{\degreeCelsius} and
    written to \plcaddr{MD48} (\code{Tin1\_HE1}).

  \item \textbf{Valve command computation.} $F_{1,SP}$ (m$^3$/h, stored at \plcaddr{MD24})
    is converted to a Word output for \plcaddr{QW8} via the inverse of the scaling formula:
    $\text{cmd} = \lfloor(F_{1,SP} / F_{1,\max}) \times 27648\rfloor$.

  \item \textbf{DB1 update FC.} A small Function copies key \plcaddr{MD} values into DB1
    at fixed offsets, making them readable by Snap7 from the PC.
\end{enumerate}

\subsection{DB1 for Snap7 Access}
\label{sec:db1}

DB1 is a Data Block configured with \textbf{standard (non-optimized) block access}.
This is a critical configuration point: by default, TIA Portal creates new DBs with
\textit{Optimized block access} enabled, which allows the compiler to arrange variables
freely for performance and hides the byte offsets. The Snap7 library requires
\textbf{known, fixed byte offsets} to read and write DB variables. Disabling optimized
access (Properties $\rightarrow$ Attributes $\rightarrow$ uncheck \textit{Optimized block
access}) makes the offsets visible in the DB editor and enables Snap7 access.

Table~\ref{tab:db1} shows the DB1 structure used by Snap7.

\begin{table}[H]
  \centering
  \caption{DB1 structure for Snap7 access (non-optimized, absolute offsets).}
  \label{tab:db1}
  \begin{tabular}{@{}C{2.2cm}C{2.4cm}L{3.5cm}L{5cm}@{}}
    \toprule
    \textbf{Byte offset} & \textbf{TIA type} & \textbf{Variable} & \textbf{Description} \\
    \midrule
    DBD0  & REAL & \code{temperature}   & $T_{in1,HE1}$ (\si{\degreeCelsius}) \\
    DBD4  & REAL & \code{flow\_rate}    & $F_1$ (\si{\metre\cubed\per\hour}) \\
    DBD8  & REAL & \code{setpoint\_temp}& Temperature setpoint (\si{\degreeCelsius}) \\
    DBD12 & REAL & \code{setpoint\_flow}& $F_{1,SP}$ (\si{\metre\cubed\per\hour}) \\
    DBW16 & INT  & \code{valve\_state}  & 0\,=\,CLOSED, 1\,=\,OPEN, 2\,=\,PARTIAL \\
    \bottomrule
  \end{tabular}
\end{table}

\section{WinCC Unified and SIMATIC Panel 7\textquotedblright}

\subsection{WinCC Unified Runtime}

WinCC Unified is Siemens' latest-generation HMI software, designed for touchscreen panels
and web browsers. The SIMATIC Unified 7\textquotedblright{} panel (IP~192.168.1.200) runs
WinCC Unified Runtime, which connects to the S7-1500 over PROFINET and maintains a local
tag cache mirroring all PLC variables.

The WinCC Unified tag table uses the same symbolic names as the PLC tag table. At runtime,
WinCC continuously polls the PLC and updates its local cache. This cache is what the
Open Pipe mechanism exposes to external applications (such as the Docker container).

\subsection{Open Pipe Communication}
\label{sec:openpipe}

Open Pipe is a Siemens proprietary mechanism that allows external processes running on the
same hardware as WinCC Unified to read and write WinCC tags via a Unix domain socket.

\textbf{Socket location.}
The socket is located at \code{/tmp/siemens/automation/openpipe.sock} on the panel
filesystem. When the Flask application runs in a Docker container on the panel, this socket
is made accessible inside the container via a volume mount:

\begin{lstlisting}[style=yaml, caption={Docker volume mount for Open Pipe access.}]
volumes:
  - /tmp/siemens/automation:/tempcontainer/
\end{lstlisting}

Inside the container, the socket appears at \code{/tempcontainer/openpipe.sock}.

\textbf{Protocol.}
Communication uses newline-delimited JSON messages over a Unix \texttt{SOCK\_STREAM} socket:

\begin{lstlisting}[style=json,
  caption={Open Pipe read/write request--response format.},
  label={lst:openpipe}]
// Read request
{"action": "read", "tags": ["Tin1_HE1", "F1_SP", "Tout2_HE1"]}
// Response
{"Tin1_HE1": 75.3, "F1_SP": 3.0, "Tout2_HE1": 42.1}

// Write request
{"action": "write", "tag": "F1_SP", "value": 3.5}
\end{lstlisting}

\textbf{Advantages over Snap7.}
Open Pipe uses tag names directly (matching the WinCC tag table), rather than raw byte
offsets. It requires no knowledge of DB structure or data types~--- WinCC handles type
conversion. Each request is served from the WinCC tag cache, without triggering a PLC poll.

\textbf{Simulation fallback.}
When the socket file does not exist (e.g., during PC development), the
\code{OpenPipeConnector} class automatically returns plausible simulated values for all
21 tags. The \code{source} field in the API response switches from \texttt{"OPENPIPE"} to
\texttt{"SIMULATION"}, alerting the user to the fallback mode.

\subsection{Snap7 and Open Pipe: Coexistence}

The project uses both protocols simultaneously but for different purposes:

\begin{table}[H]
  \centering
  \caption{Coexistence of Snap7 and Open Pipe communication layers.}
  \label{tab:protocols}
  \begin{tabular}{@{}L{2.5cm}L{3.5cm}C{1.8cm}L{4cm}@{}}
    \toprule
    \textbf{Protocol} & \textbf{Module} & \textbf{Tags} & \textbf{Used by} \\
    \midrule
    Snap7 (TCP 102)   & \code{plc\_connector.py}     & 5 (DB1)  & Monitor page, CSV logger \\
    Open Pipe (socket)& \code{openpipe\_connector.py}& 21 (WinCC)& Synoptique, Cascade, Process \\
    \bottomrule
  \end{tabular}
\end{table}

The two connectors operate on independent threads and do not interfere with each other.

\section{Industrial Edge Deployment}

\subsection{Industrial Edge Overview}

Siemens Industrial Edge turns compatible hardware (including the SIMATIC Unified panel)
into an embedded Linux host running a Docker daemon managed by the Edge Runtime service.
The deployment workflow introduces a packaging step: the application must be bundled into
a \texttt{.app} file using the \textit{Siemens Industrial Edge App Publisher} desktop tool.
This tool reads \code{docker-compose.edge.yml}, embeds the Docker image layers, and
produces a signed archive for the Edge Management console.

\subsection{Edge Docker Compose Configuration}

The edge-specific \code{docker-compose.edge.yml} differs from the development configuration
in three key respects:

\begin{lstlisting}[style=yaml,
  caption={Edge Docker Compose configuration (\texttt{docker-compose.edge.yml}).},
  label={lst:edgecompose}]
services:
  app:
    image: project-s7-app:latest
    restart: unless-stopped
    ports:
      - "30500:5000"      # Edge apps must use ports 30000-35000
    environment:
      - FLASK_DEBUG=0
      - LOG_DIR=/app/logs
    mem_limit: 256m       # Panel hardware constraint
    volumes:
      - /tmp/siemens/automation:/tempcontainer/
      - ie-databus:/app/logs
    networks:
      - ieguestnetwork
\end{lstlisting}

\begin{itemize}[leftmargin=*, itemsep=3pt]
  \item \textbf{Port 30500:} Industrial Edge requires user-application ports in the range
    30\,000--35\,000.
  \item \textbf{Memory limit:} 256\,MB respects the panel's limited hardware resources.
  \item \textbf{Volume mount:} \code{/tmp/siemens/automation} grants the container access
    to the Open Pipe socket.
  \item \textbf{Guest network:} \code{ieguestnetwork} provides connectivity to the panel's
    PROFINET interface without exposing internal services.
\end{itemize}

\subsection{Deployment Process and Current Status}

The standard deployment sequence is:
\begin{enumerate}[leftmargin=*, itemsep=3pt]
  \item Build the Docker image on the development PC:
    \code{docker compose -f docker-compose.edge.yml build}
  \item Import \code{docker-compose.edge.yml} into App Publisher $\rightarrow$
    validate $\rightarrow$ export as \texttt{.app}
  \item Log in to the Edge Management console at \code{http://192.168.1.200} with
    Industrial Edge licence credentials $\rightarrow$ upload \texttt{.app} $\rightarrow$ deploy
  \item Access the deployed application at \code{http://192.168.1.200:30500}
\end{enumerate}

\textbf{Current status:} The \texttt{.app} file has been successfully generated by
App Publisher. Deployment on the panel is pending acquisition of Industrial Edge licence
credentials. In the interim, the application runs on the development PC (port 5000) and is
accessible from the panel browser at \code{http://192.168.1.149:5000}. A Windows Firewall
inbound rule was added to allow TCP connections on port~5000 from the local subnet.
